% Generated draft: requires mathematical and source review.
\documentclass[11pt]{article}
\usepackage[margin=25mm]{geometry}
\usepackage{fontspec}
\setmainfont{Latin Modern Roman}
\usepackage{amsmath,amssymb,mathtools,unicode-math}
\setmathfont{Latin Modern Math}
\usepackage{enumitem,needspace,xurl,hyperref}
\setlength{\parindent}{0pt}
\setlength{\parskip}{5pt}
\setlength{\emergencystretch}{3em}
\newcommand{\N}{\mathbb{N}}
\newcommand{\Z}{\mathbb{Z}}
\newcommand{\Q}{\mathbb{Q}}
\newcommand{\R}{\mathbb{R}}
\newcommand{\C}{\mathbb{C}}
\newcommand{\F}{\mathbb{F}}
\newcommand{\A}{\mathbb{A}}
\newcommand{\PP}{\mathbb{P}}
\newcommand{\E}{\mathbb{E}}
\newcommand{\eps}{\varepsilon}
\newcommand{\dd}{\,\mathrm{d}}
\begin{document}
\section*{UnsolvedMath ID: 11000086}
\textbf{Source problem code:} catalogue entry
\subsection{Definitions and mathematical statement}
Let \(H\leq\operatorname{Mod}_g\). The source asks to compute two distinct densities: that of \(\operatorname{spec}(H)\) inside \(\operatorname{spec}(\operatorname{Mod}_g)\), and that of \(H\cap P_g\) inside \(P_g\). It specifically asks for the spectral density of the level subgroup \(\operatorname{Mod}_g[L]\), and for the analogous questions with \(H=\mathcal I_g(k)\), \(k\geq1\).

The preceding definitions of \(\operatorname{spec}\), \(P_g\) and density are missing. In the pseudo-Anosov context, a possible spectral convention is the set of stretch factors (or their logarithms) of pseudo-Anosov elements of a subgroup. Once a locally finite height-counting convention is fixed, a distinct-value spectral density could be
\[\lim_{R\to\infty}\frac{\#\{a\in\operatorname{spec}(H):a\leq R\}}
{\#\{a\in\operatorname{spec}(\operatorname{Mod}_g):a\leq R\}}.\]
Multiplicity-weighted spectra or counts of conjugacy classes instead define different quantities. If \(P_g\) denotes pseudo-Anosov elements and a finite generating set \(S\) supplies the sample, a possible element-density convention would be
\[\lim_{R\to\infty}\frac{|H\cap P_g\cap B_S(R)|}{|P_g\cap B_S(R)|}.\]
These displays are candidate specifications, not definitions inferred from the source; one must state the exhaustion/measure, multiplicities and whether limits, upper/lower densities or another notion are intended. Counting all elements by a conjugacy-invariant spectral cutoff can give infinite fibers, so that convention cannot be assumed without a quotient or other finite sampling rule.

In standard notation, the level group is the kernel of the homology action modulo \(L\), and \(\mathcal I_g\) is the kernel of the integral homology action. The \(\mathcal I_g(k)\) filtration is generally defined through actions on nilpotent quotients of \(\pi_1\Sigma_g\), but its precise indexing and outer-action convention are not given here. The genus/level ranges are likewise not repeated. The request is to compute densities for relevant subgroups, not an assertion that either density always exists or that spectral density equals element density.
\subsection{Short English statement}
Measure how often subgroup pseudo-Anosov data occur in the full mapping class group, both as spectral values and as elements, especially for level and Torelli-filtration groups. The spectrum and density sampling conventions are missing.
\subsection{Sources}
\begin{itemize}
\item Local UnsolvedMath registry, permanent ID \(11000086\), exact source code \texttt{catalogue entry}.
\item Farb, \emph{Problems on Mapping Class Groups and Related Topics}, Question 7.12, PDF page 59: \url{https://www.math.uchicago.edu/~farb/papers/mcgbook.pdf}. Registry access 2026-07-29; pdf-font-structure extraction. No separate density/result reference is supplied.
\end{itemize}
\end{document}
